% संपूर्ण मराठी ग्रंथातील प्रकरण 34: द्वितीय-क्रम तर्कशास्त्र आणि संचसिद्धान्त
% हा संपादनयोग्य स्रोतखंड आहे. संकलनासाठी संपूर्ण स्रोतसंग्रहातील
% अक्षररूपे, पूर्वभाग, चिन्हव्याख्या आणि आकृती-स्रोत आवश्यक आहेत.
% मूळ: Open Logic Project; मजकूर आणि रूपांतर CC BY 4.0.
% यंत्रानुवाद, दुरुस्ती आणि तपासणी: OpenAI Codex —
% GPT-5.6 Sol आणि GPT-6 Sol, दोन्ही Ultra effort.
% दुरुस्ती-पुनरावलोकन: GPT-6.1 Sol, Ultra effort.
\chapter{द्वितीय-क्रम तर्कशास्त्र आणि संचसिद्धान्त}\label{sol:set:chap}
\begin{editorial}
या विभागात द्वितीय-क्रम तर्कशास्त्रात घातसंच आणि
सांतत्य यांचे संकेतीकरण केले आहे. निष्कर्ष मांडले
आहेत; पण त्यांचे पुरावे अजून भरायचे आहेत. अद्याप
सरावप्रश्न नाहीत---आणि व्याख्या व निष्कर्षांमध्येही
त्रुटी असू शकतात. सावधपणे वापरा; काही असत्य
किंवा अस्पष्ट आढळल्यास कळवा.
\end{editorial}









\section{परिचय}\label{sol:set:int:sec}

द्वितीय-क्रम तर्कशास्त्रात प्रांताच्या उपसंचांवर तसेच
फलनांवर संख्यापन करता येते; म्हणून त्यात संचसिद्धान्ताचा
काही भाग तरी मांडता येईल अशी अपेक्षा आहे. येथे
‘मांडणे’ याचा अर्थ संचसिद्धान्तीय गुणधर्म आणि विधाने
द्वितीय-क्रम तर्कशास्त्रात व्यक्त करता येणे, तेही
संचांसाठी कोणतीही खास अतार्किक चिन्हे (उदा.,
संचसिद्धान्तातील सदस्यत्वाचे विधेय) न वापरता.
उदाहरणार्थ, संचांचे संयोग आणि छेद तसेच उपसंच-संबंध
परिभाषित करता येतात; शिवाय संचांच्या आकारांची तुलना
करता येते आणि कँटरच्या प्रमेयासारखे निष्कर्षही मांडता येतात.












\section{संचांची तुलना}\label{sol:set:cmp:sec}

\begin{prop}
सूत्र $\lforall[x][(X(x) \lif Y(x))]$ उपसंच-संबंध
परिभाषित करते; म्हणजे
$\Sat{M}{\lforall[x][(X(x) \lif Y(x))]}[s]$
जर आणि फक्त जर $s(X) \subseteq s(Y)$.
\end{prop}

\begin{prop}
सूत्र $\lforall[x][(X(x) \liff Y(x))]$ संचांवरील
समानता-संबंध परिभाषित करते; म्हणजे
$\Sat{M}{\lforall[x][(X(x) \liff Y(x))]}[s]$
जर आणि फक्त जर $s(X) = s(Y)$.
\end{prop}

\begin{prop}
सूत्र $\lexists[x][X(x)]$ अरिक्त असण्याचा
गुणधर्म परिभाषित करते; म्हणजे
$\Sat{M}{\lexists[x][X(x)]}[s]$
जर आणि फक्त जर $s(X) \neq \emptyset$.
\end{prop}

संच~$X$ हा संच~$Y$ पेक्षा आकारमानाने मोठा नाही,
$\cardle{X}{Y}$, जर आणि फक्त जर
एकास-एक फलन $f\colon X \to Y$ असेल.
फलन एकास-एक आहे, तसेच~$X$ मधील निविष्टांसाठी
त्याची मूल्ये~$Y$ मध्ये आहेत, हे व्यक्त करता येते.
म्हणून प्रांताच्या उपसंचांवर आकारमानाने मोठा नाही
हा संबंधही परिभाषित करता येतो.

\begin{prop}

पुढील सूत्र
\[\lexists[u][(\lforall[x][(X(x) \lif Y(u(x)))] \land \lforall[x][\lforall[y][((X(x) \land X(y) \land \eq[u(x)][u(y)]) \lif
      \eq[x][y])]])]
\]
आकारमानाने मोठा नाही हा संबंध परिभाषित करते. येथे
एकास एक अट फक्त X मधील दोन घटकांसाठी मागितली आहे.
X वरील फलन पूर्ण प्रांतावर मनमानी वाढवता येते.
\end{prop}

दोन संचांचे आकारमान समान असते, म्हणजे ते
``तुल्यबल'' असतात, $\cardeq{X}{Y}$, जर आणि
फक्त जर एकास-एक व आच्छादक फलन~$f\colon X \to Y$ असेल.

\begin{prop}
पुढील सूत्र
\begin{multline*}
  \lexists[u][(\lforall[x][(X(x) \lif Y(u(x)))] \land {}\\
    \lforall[x][\lforall[y][((X(x) \land X(y) \land \eq[u(x)][u(y)]) \lif \eq[x][y])]]
  \land {} \\\lforall[y][(Y(y) \lif \lexists[x][(X(x)
      \land \eq[y][u(x)])])])]
\end{multline*}
तुल्यबल असण्याचा संबंध परिभाषित करते.
\end{prop}

या सूत्रे ना अनुक्रमे $X \subseteq Y$,
$X = Y$, $X \neq \emptyset$, $\cardle{X}{Y}$
आणि $\cardeq{X}{Y}$ अशी संक्षिप्त चिन्हे वापरू.
(संच $X$ आणि $Y$ यांच्याविषयी अनौपचारिकपणे
बोलतानाही हीच चिन्हे वापरल्यामुळे थोडा गोंधळ होऊ शकतो;
पण येथे ती एकस्थानी संबंधचले $X$ आणि~$Y$
असलेल्या द्वितीय-क्रम तर्कशास्त्रातील सूत्रे
साठीची संक्षिप्त चिन्हे आहेत.)

\begin{prop}
वाक्य $\lforall[X][\lforall[Y][((\cardle{X}{Y} \land
    \cardle{Y}{X}) \lif \cardeq{X}{Y})]]$ वैध आहे.
\end{prop}

\begin{proof}
कोणत्याही उपसंचांसाठी $X \subseteq \Domain{M}$
आणि $Y \subseteq \Domain{M}$, जर
$\cardle{X}{Y}$ आणि $\cardle{Y}{X}$, तर
$\cardeq{X}{Y}$ असेल, तर संरचना~$\Struct{M}$
मध्ये या वाक्याची पूर्ति होते. पण हे
\emph{कोणत्याही} संच $X$ आणि $Y$ साठी खरे आहे:
हेच श्रेडर--बर्नस्टाइन प्रमेय आहे.
\end{proof}













\section{संचांची आकारमाने}\label{sol:set:crd:sec}

\begin{explain}
प्रांत सांत किंवा अनंत, गणनीय किंवा
अगणनीय आहे हे जसे व्यक्त करता येते,
तसेच~$\Domain{M}$ चा उपसंच सांत किंवा अनंत,
गणनीय किंवा अगणनीय असण्याचा
गुणधर्मही परिभाषित करता येतो.
\end{explain}

\begin{prop}

सूत्र $\fn{Inf}(X) \ident$
\begin{multline*}
\lexists[u][(\lforall[x][(X(x) \lif X(u(x)))] \land {} \\
\lforall[x][\lforall[y][((X(x) \land X(y) \land \eq[u(x)][u(y)]) \lif
      \eq[x][y])]] \land {} \\
\lexists[y][(X(y) \land \lforall[x][(X(x) \lif \eq/[y][u(x)])])])]
\end{multline*}
याची चल-मूल्यांकन~$s$ नुसार पूर्ति होईल, जर आणि
फक्त जर $s(X)$ अनंत असेल. स्रोतदुरुस्ती: फलन
X ला स्वतःमध्येच पाठवते, X वर एकास एक असते आणि X मधील
किमान एक घटक त्याच्या X वरील व्याप्तीत नसतो, अशा सर्व अटी
घेतल्या आहेत. त्यामुळे मूळ सूत्राचा सांत उपसंचावरील counterexample
दूर होतो. अनंततेचा नेहमीचा निवडीसह संदर्भ येथेही लागू आहे.
\end{prop}

\begin{prop}
सूत्र $\fn{Count}(X) \ident $
\begin{multline*}
\lnot\lexists[z][X(z)] \lor {} \\
\lexists[z][\lexists[u][(X(z) \land
    \lforall[x][(X(x) \lif X(u(x)))] \land {} \\
\lforall[Y][((Y \subseteq X \land Y(z) \land
      \lforall[x][(Y(x) \lif Y(u(x)))]) \lif X = Y)])]]
\end{multline*}
याची चल-मूल्यांकन~$s$ नुसार पूर्ति होईल, जर आणि
फक्त जर $s(X)$ गणनीय असेल. स्रोतदुरुस्ती:
पहिला खंड रिकामा संच घेतो. अरिक्त X साठी आरंभीचा घटक
आणि X ला स्वतःमध्ये पाठवणारे फलन घेतले आहे. शेवटची अट
फक्त X मधील बंद उपसंचांबद्दल आहे; त्यामुळे X हा आरंभीच्या
घटकाचा पूर्ण orbit असतो. हे संपूर्ण प्रांतच X असण्याची
चुकीची सक्ती करत नाही.
\end{prop}

कँटरच्या प्रमेयावरून अगणनीय संच आहेत,
आणि अनंत आकारमानांचे निरनिराळे अमर्याद स्तरही
आहेत, हे आपल्याला माहीत आहे. संचसिद्धान्तात
संचांच्या आकारमानांचे संपूर्ण अंकगणित विकसित केले
जाते आणि संचांना अनंत संचांक दिले जातात. सांत
संचांची आकारमाने मोजणारे संचांक म्हणजे स्वाभाविक
संख्या. गणनीय अनंत संचांचा संचांक हा पहिला
अनंत संचांक असून त्याला~$\aleph_0$ (``अलेफ-शून्य'')
म्हणतात. पुढील अनंत आकारमान~$\aleph_1$ आहे.
गणनीय नसलेल्या संचाला असू शकणारे हे
सर्वांत लहान आकारमान आहे (गणनीय अनंत संचाचे
आकारमान~$\aleph_0$ असते). ``$X$ चे आकारमान
$\aleph_0$ आहे'' याची व्याख्या
$\fn{Aleph}_0(X) \liff \fn{Inf}(X) \land \fn{Count}(X)$
अशी करता येते. $X$ चे आकारमान $\aleph_1$
असेल, जर आणि फक्त जर त्याच्यापेक्षा काटेकोरपणे लहान
आकारमानाचे सर्व उपसंच सांत किंवा~$\aleph_0$ आकारमानाचे
असतील आणि तो स्वतः सांत किंवा~$\aleph_0$ आकारमानाचा
नसेल. म्हणून पुढील दुरुस्त सूत्र घेतो:
$\fn{Aleph_1}(X) \ident \lnot\fn{Count}(X) \land
\lforall[Y][((Y \subseteq X \land \lnot\cardeq{Y}{X}) \lif
  (\lnot\fn{Inf}(Y) \lor \fn{Aleph}_0(Y)))]$. $\aleph_2$ आकारमानाची
व्याख्याही अशाच प्रकारे करायची, असे त्यात म्हटले आहे.
स्रोतदुरुस्ती: X अगणनीय असावा लागतो. स्वतः X किंवा
त्याच्याशी तुल्यबल असलेला उपसंच हा लहान आकारमानाचा नाही;
त्याला finite किंवा aleph-zero असण्याची चुकीची अट लावलेली नाही.

एक विशेष महत्त्वाचे आकारमान म्हणजे तथाकथित
सांतत्याचा संचांक. ते $\Pow{\Nat}$ चे आकारमान,
किंवा समतुल्यपणे~$\Real$ चे आकारमान आहे.
एखादा संच सांतत्याच्या आकारमानाचा आहे, हेही
द्वितीय-क्रम तर्कशास्त्रात व्यक्त करता येते; मात्र
त्यासाठी आणखी थोडे काम करावे लागते.












\section{सांतत्याचे आकारमान}\label{sol:set:pow:sec}

\begin{explain}

द्वितीय-क्रम तर्कशास्त्रात प्रांताच्या उपसंचांवर
संख्यापन करता येते; पण प्रांताच्या उपसंचांच्या
संचांवर करता येत नाही. ते थेट करण्यासाठी
\emph{तृतीय-क्रम} तर्कशास्त्र लागेल. उदाहरणार्थ,
संचाच्या घातसंचापासून मूळ संचाकडे एकास-एक
फलन नसते हे कँटरचे प्रमेय मांडायचे असल्यास आपण
असे म्हणण्याचा प्रयत्न करू शकतो: ``प्रत्येक
संच~$X$ आणि प्रत्येक संच~$P$ यांसाठी, $P$
हा~$X$ चा घातसंच असेल, तर
$\cardle{P}{X}$ नसते.'' पण $P$ हा~$X$ चा
घातसंच आहे असे म्हणण्यासाठी $P$ चे घटक
नेमके~$X$ चे सर्व उपसंच आहेत, हे सूत्रबद्ध करावे
लागेल; उदाहरणार्थ,
$\lforall[Y][(P(Y) \liff Y \subseteq X)]$.
अडचण $P(Y)$ मध्ये आहे: ते द्वितीय-क्रम
तर्कशास्त्रातील सूत्र नाही, कारण
$P$ सारख्या एकस्थानी संबंधचलाचे निविष्ट फक्त
पदेच असू शकतात.

मात्र प्रांत पुरेसा मोठा असेल, तर संचांच्या संचांवरचे
संख्यापन \emph{अनुकरणाने} करता येते. कल्पना अशी
की द्विस्थानी संबंध~$R$ हा प्रांत मधील
घटक ना त्याच प्रांत मधील
घटक शी जोडतो. असा~$R$ दिल्यास,
एखाद्या~$x$ शी $R$-संबंधित सर्व घटक
गोळा करता येतात: $\Setabs{y \in
  \Domain{M}}{R(x,y)}$ हा~$x$ ने
``सांकेतिक केलेला'' संच आहे. उलट,
$Z \subseteq \Pow{\Domain{M}}$ हा~$\Domain{M}$
च्या उपसंचांचा संग्रह असेल आणि~$\Domain{M}$
मध्ये~$Z$ मधील संचांइतके तरी घटक
असतील, तर असा संबंध~$R \subseteq \Domain{M}^2$
असतो की प्रत्येक $Y \in Z$ हा
काही~$x$ ने~$R$ द्वारे सांकेतिक केला जातो.
\end{explain}

\begin{defn}
$R \subseteq \Domain{M}^2$ असेल, तर
$\Setabs{y \in \Domain{M}}{R(x, y)}$ या संचाचे
\emph{$x$ ने $R$-सांकेतीकरण} होते.
\end{defn}

घटक~$x \in \Domain{M}$ हा
$Z \subseteq \Domain{M}$ संचाचे $R$-सांकेतीकरण
करत असेल, तर संच $Y \subseteq \Domain{M}$ हा
संचांच्या संचाचे सांकेतीकरण करतो: ते म्हणजे~$Y$
मधील घटक नी सांकेतिक केलेले संच.
म्हणून संच $Y$ हा $\Pow{X}$ चे
$R$-सांकेतीकरण करू शकतो. असे होईल, जर
आणि फक्त जर प्रत्येक $Z \subseteq X$ साठी
काही $x \in Y$ हे~$Z$ चे $R$-सांकेतीकरण
करत असेल आणि प्रत्येक $x \in Y$ हा
काही~$Z \subseteq X$ चे $R$-सांकेतीकरण करत असेल.

\begin{prop}
सूत्र
  \[  \fn{Codes}(x, R, Z) \ident \lforall[y][(Z(y) \liff
    R(x, y))]
  \]
हे $s(x)$ हा $s(R)$ द्वारे $s(Z)$ चे
सांकेतीकरण करतो, असे व्यक्त करते. पुढील
सूत्र
\begin{multline*}
  \fn{Pow}(Y, R, X) \ident \\
  \lforall[Z][(Z \subseteq X \lif \lexists[x][(Y(x) \land
    \fn{Codes}(x, R, Z))])] \land {} \\ \lforall[x][(Y(x) \lif
  \lforall[Z][(\fn{Codes}(x, R, Z) \lif Z \subseteq X)])]
\end{multline*}
हे $s(Y)$ हा $s(R)$ द्वारे~$s(X)$ च्या
घातसंचाचे सांकेतीकरण करतो, असे व्यक्त करते;
म्हणजे $s(Y)$ चे घटक $s(R)$ द्वारे
नेमके $s(X)$ चे सर्व उपसंच सांकेतिक करतात.
\end{prop}

\begin{explain}
या युक्तीने उपसंचांवर थेट संख्यापन करण्याऐवजी
त्यांच्या संकेतांवर संख्यापन करून घातसंचाविषयी
विधाने व्यक्त करता येतात. उदाहरणार्थ,
आता कँटरचे प्रमेय असे मांडता येते की~$X$ चा
घातसंच सांकेतिक करणाऱ्या कोणत्याही संबंधाच्या
प्रांतापासून मूळ~$X$ कडे एकास-एक फलन नसते.
\end{explain}

\begin{prop}
पुढील वाक्य
\begin{multline*}
  \lforall[X][\lforall[Y][\lforall[R][(\fn{Pow}(Y, R, X) \lif \\
      \lnot \lexists[u][(\lforall[x][\lforall[y][((Y(x) \land Y(y) \land \eq[u(x)][u(y)]) \lif
            \eq[x][y])]] \land {}\\
        \lforall[x][(Y(x) \lif X(u(x)))])])]]]
\end{multline*}
वैध आहे.
\end{prop}

\begin{explain}
गणनीय अनंत संचाचा घातसंच अगणनीय
असतो; म्हणून त्याचा संचांक कोणत्याही
गणनीय अनंत संचाच्या संचांकापेक्षा
मोठा असतो (तो संचांक $\aleph_0$ असतो).
$\Pow{\Nat}$ चे आकारमान ``सांतत्याचे आकारमान''
म्हणतात, कारण ते वास्तव संख्यारेषेवरील
बिंदूंच्या संचाच्या, म्हणजे~$\Real$ च्या,
आकारमानाइतके असते. एखाद्या
गणनीय अनंत संचाचा घातसंच सांकेतिक
करण्याइतका प्रांत मोठा असेल, तर एखादा
संच सांतत्याच्या आकारमानाचा आहे हे असे
मांडता येते: तो~$\aleph_0$ आकारमानाच्या संच~$X$
चा घातसंच सांकेतिक करणाऱ्या कोणत्याही
संच~$Y$ शी तुल्यबल आहे. (प्रांत पुरेसा मोठा
नसेल, म्हणजे त्यात~$\Real$ शी तुल्यबल
उपसंचच नसेल, तर~$\Pow{X}$ चे सांकेतीकरण
करणारा संबंधही असू शकत नाही.)
\end{explain}

\begin{prop}
$\cardle{\Real}{\Domain{M}}$ असेल, तर पुढील
सूत्र
\begin{multline*}
\fn{Cont}(Y) \ident
\lexists[X][\lexists[R][((\fn{Aleph}_0(X) \land
      \fn{Pow}(Y, R, X)) \land \\
      \lforall[x][\lforall[y][((Y(x) \land {}
      Y(y) \land \lforall[z][(R(x,z) \liff R(y,z))]) \lif \eq[x][y])]])]]
\end{multline*}
हे $\cardeq{s(Y)}{\Real}$ असे व्यक्त करते.
\end{prop}

\begin{proof}
  $\fn{Pow}(Y, R, X)$ यानुसार $s(Y)$ हा $s(R)$
  द्वारे~$s(X)$ चा घातसंच सांकेतिक करतो;
  $\fn{Aleph}_0(X)$ नुसार हा संच गणनीय अनंत असतो.
  म्हणून $s(Y)$ चे आकारमान किमान सांतत्याच्या
  आकारमानाइतके असते; ते अधिकही असू शकते
  (जर $s(Y)$ चे अनेक घटक~$X$ चा तोच उपसंच
  सांकेतिक करत असतील). शेवटचे संयुक्तपद
  हे टाळते: $s(R)$ द्वारे $s(Y)$ चे घटक आणि
  $s(X)$ चे उपसंच यांच्यातील संगती
  एकास-एक असावी, अशी त्याची अट आहे.
  उलट, Y सांतत्याच्या आकारमानाचा असेल, तर प्रांतातील
  गणनीय अनंत उपसंच X च्या घातसंचाशी त्याची bijection निवडता
  येते. त्या bijection मधील सदस्यत्वाने R परिभाषित केल्यास
  Pow आणि unique-code या दोन्ही अटी मिळतात. येथे आधीच्या
  विभागातील दुरुस्त Count आणि Aleph-zero definitions वापरल्या आहेत.
\end{proof}

\begin{prop}
$\cardeq{\Domain{M}}{\Real}$ जर आणि फक्त जर
\begin{multline*}
  \Sat{M}{\lexists[Y][(\lforall[x][Y(x)] \land \fn{Cont}(Y))]}.
        \end{multline*}
  स्रोतदुरुस्ती: Y हाच पूर्ण प्रांत असणे आणि त्याचे
  आकारमान Cont ने ओळखले जाणे स्पष्ट मागितले आहे. फक्त
  काही घातसंच code करणे व जुन्या range अटी घेणे अपुरे होते:
  identity फलनाने त्या अटी अधिक मोठ्या प्रांतातही पूर्ण होतात.
  \end{prop}

\begin{explain}
सांतत्य परिकल्पनेनुसार सांतत्याचे आकारमान
हा पहिला अगणनीय संचांक असतो;
म्हणजे $\Pow{\Nat}$ चे आकारमान~$\aleph_1$ असते.
\end{explain}

\begin{prop}
सांतत्य परिकल्पना खरी असेल, जर आणि फक्त जर
\[\fn{CH} \ident
\lforall[X][(\fn{Aleph}_1(X) \liff \fn{Cont}(X))]\]
वैध असेल. येथे आधी दुरुस्त केलेल्या आकारमानांच्या
definitions वापरल्या आहेत. सांतत्य आणि अलेफ-एक समान असतील,
तर दोन्ही गुणधर्म प्रत्येक उपसंचासाठी समान ठरतात. उलट,
सांतत्याच्या आकारमानाचा प्रांत घेऊन त्यातील अलेफ-एक
आकारमानाच्या उपसंचाला हे विधान लावल्यास अपेक्षित समानता मिळते.
\end{prop}

हे लक्षात घ्या की $\lnot \fn{CH}$ वैध असेल,
जर आणि फक्त जर सांतत्य परिकल्पना खोटी असेल,
असे होत नाही. गणनीय प्रांतात
आकारमान~$\aleph_1$ असलेले उपसंच नसतात,
आणि सांतत्याच्या आकारमानाचे उपसंचही नसतात;
म्हणून अशा गणनीय प्रांतात
$\fn{CH}$ नेहमी खरे असते. मात्र सांतत्य
परिकल्पना खोटी असेल, जर आणि फक्त जर
वैध असेल असे पुढील वेगळे वाक्य देता येते:

\begin{prop}
सांतत्य परिकल्पना खोटी असेल, जर आणि
फक्त जर \[\fn{NCH} \ident
\lforall[X][(\fn{Cont}(X) \lif \lexists[Y][(Y \subseteq X \land \lnot
    \fn{Count}(Y) \land \lnot \cardeq{X}{Y})])]\]
वैध असेल. सांतत्य अलेफ-एकपेक्षा मोठे असेल, तर
प्रत्येक सांतत्य आकारमानाच्या संचात अलेफ-एक आकारमानाचा
अगणनीय, त्याच्याशी तुल्यबल नसलेला उपसंच निवडता येतो.
सांतत्य अलेफ-एक असेल, तर अशा लहान आकारमानाचा अगणनीय
उपसंच नसतो. लहान प्रांतांत Cont-संच नसल्याने शर्त रिक्ततः खरी असते.
\end{prop}



\part{लॅम्डा कलन}\label{lam:part}
\begin{editorial}
हा भाग लॅम्डा कलनाविषयी आहे. परिचयाचे प्रकरण जेरेमी अविगाड यांच्या
टिपणांवर आधारित आहे; त्यातील काही भाग आता पुनरुक्त झाला असून पुढील
प्रकरणांमध्ये समाविष्ट आहे. विन्यासमीमांसा, चर्च--रॉसर गुणधर्म आणि
लॅम्डा-परिभाष्यता यांवरील प्रकरणे झेसन चियान यांनी मिटॅक्सच्या उन्हाळी
प्रशिक्षणकाळात तयार केली. त्यांचे अजून परीक्षण आणि पुनरीक्षण करायचे आहे.
\end{editorial}
